% !TEX TS-program = lualatex
% Bridge tkz-fn -> tkz-elements (points exported to z, in cm)
\documentclass{article}
\usepackage[margin=2cm]{geometry}
\usepackage{tkz-elements}
\usepackage[fr]{tkz-fn}
\parindent=0pt
\begin{document}

Repère non orthonormé : xstep=1, ystep=2 (1 cm représente 2 unités en ordonnée).

\tkzFnDef{f}{x^2/2}
\begin{tikzpicture}
  \tkzFnInit[xmin=-3,xmax=3,ymin=-2,ymax=6,ystep=2]
  \tkzFnGrid
  \tkzFnAxes
  \tkzFnDraw[thick,red](f)
  \tkzFnDefPoint(f,2){A}       % A(2;2) in the frame, z.A = (2,1) in cm
  \tkzFnDefPointXY(0,0){O}
  \directlua{
    C.OA = circle(z.O, z.A)     % tkz-elements works in cm
    z.B  = C.OA.south
  }
  \tkzGetNodes
  \draw[blue] (O) circle [radius=\tkzUseLua{C.OA.radius}];
  \tkzFnDrawPoints[red](A,O)
  \fill[blue] (B) circle (1.5pt) node[below right] {$B$};
  \node[above left] at (A) {$A$};
\end{tikzpicture}

$A\tkzFnPrintCoord(A)$ dans le repère ; dans tkz-elements :
z.A = (\directlua{tex.sprint(z.A.re)}, \directlua{tex.sprint(z.A.im)}) en cm.

Retour vers le repère : B a pour coordonnées
\directlua{tkzfn.point_from_z("B")}%
$\tkzFnPrintCoord(B)$ (rayon $\sqrt5$ cm donc $y_B=-2\sqrt5\approx -4{,}47$).
\end{document}
